\documentclass{resume}

\usepackage[hidelinks]{hyperref}

\begin{document}
\pagenumbering{gobble}

\noindent
\name{李昱辰}
\otherInfo{27 届}{意向岗位：Python Agent / LLM 应用开发实习}{}{}
\otherInfo{可尽快到岗}{每周可实习 5 天}{可连续实习 6 个月}{}
\otherInfo{邮箱：edllyc041202@gmail.com}{GitHub：\href{https://github.com/EDLlyc/12306-ticket-service}{EDLlyc/12306-ticket-service}}{}{}

\vspace{-1.0em}

\section{教育背景}
\datedsubsection{\textbf{安徽大学} - 数字媒体技术专业 - 本科}{2023.09 - 2027.06}
\begin{itemize}[parsep=0.5ex]
    \item \textbf{科研与竞赛成果}：以学生一作身份向 \textit{IEEE TMI} 在投算法论文 1 篇（外审），提出基于拓扑学先验的连通性损失函数；获美赛二等奖、数学建模省赛三等奖、三等学业奖学金。
    \item \textbf{外语能力}：雅思（IELTS）7.0 分，具备英文技术文档、学术论文阅读与复现能力。
    \item \textbf{成绩情况}：均分 87.6/100，系统修读数据结构、操作系统、计算机网络、数据库等计算机核心课程。
\end{itemize}

\section{实习经历}
\datedsubsection{\textbf{赛先文化} - AI Agent 应用开发实习生}{2026.07 - 2026.08}
\textbf{项目技术}：Python, FastAPI, LangGraph, PostgreSQL, pgvector, Redis, Docker, 智谱 GLM, RAG, Tool Calling\\
\textbf{业务方向}：参与教育行业新媒体数字员工建设，围绕新闻采集、事实治理、智能选题、品牌文案、配图审核与企业微信交付，搭建可追溯、可评测、可恢复的内容生产链路。\\
\textbf{主要工作}：
\begin{enumerate}[parsep=0.5ex]
    \item \textbf{Agent 工作流与智能选题}：基于 LangGraph 和持久化任务状态编排新闻采集、事件治理、选题、文案与配图流程；设计“规则阈值过滤 + LLM 受控重排”机制，模型仅对合格候选排序，并通过优先级屏障、结果校验和确定性回退避免异常输出影响业务。
    \item \textbf{证据治理与品牌 RAG}：对新闻来源建立不可变快照、事实标签、事件版本和引用绑定，保证生成内容可回溯；基于 PostgreSQL/pgvector 构建数字 IP 与品牌知识检索，将新闻事实和品牌表达分离，降低事实混淆与品牌内容误用风险。
    \item \textbf{多模态内容生产}：接入智谱文本模型和图像生成服务，完成家长友好文案、品牌配图与素材包生成；实现图片格式校验、相似度检查、可选 OCR、失败重试与安全备选图机制，避免供应商异常阻塞后续交付。
    \item \textbf{可靠交付与工程质量}：基于 FastAPI、AsyncIO 和 SQLAlchemy 实现异步服务与后台任务，完善幂等、checkpoint、超时、重试和故障恢复；建设离线评测、类型检查、契约测试及发布回退流程，并将通过审核的素材包交付至企业微信内部流程。
\end{enumerate}

\section{专业技能}
\begin{itemize}[parsep=0.5ex]
    \item \textbf{Python Agent / FastAPI / AsyncIO / HTTPX}：具备 Python 智能体服务开发经验，熟悉 FastAPI 接口设计、异步 I/O、HTTP 工具调用、SSE 流式输出与工程化测试。
    \item \textbf{LangGraph / Tool Calling / MCP / RAG / DeepEval}：熟悉 LangGraph 状态图、多意图路由、Tool Calling、MCP 外部工具接入、Plan-and-Execute/Replan、Milvus/OpenSearch 混合检索、RRF/Reranker 与 DeepEval/RAGAS 离线评测。
    \item \textbf{Java / JVM / 并发}：掌握 Java 核心基础，熟悉集合、并发编程与 JVM 内存模型，具备基于 CompletableFuture 的异步编排与高并发链路设计经验。
    \item \textbf{Spring Boot / MyBatis-Plus / MySQL}：熟练使用 Spring Boot、MyBatis-Plus、MySQL 进行后端服务开发，具备索引优化、事务隔离、SQL 调优与业务链路落库实践经验。
    \item \textbf{Redis / Redisson / Lua / RocketMQ / Sentinel / Docker / Git}：掌握 Redis/Redisson/Lua 缓存与并发控制方案，具备多级缓存、分布式锁、库存原子扣减、MQ 事务消息/延迟消息、限流熔断治理经验。
\end{itemize}

\section{项目经历}

\datedsubsection{\textbf{12306 智能数字员工}（Python Agent / RAG 智能体研发）}{2026.03 - 2026.06}
\textbf{项目技术}：Python, FastAPI, LangGraph, LangChain, MCP, Redis, Milvus, OpenSearch, 智谱 GLM, DeepEval, RAGAS\\
\textbf{项目描述}：面向 12306 规章问答与购退改查场景，将智能体编排层重构为 Python 服务；Java 票务系统作为高并发交易工具内核，由 Python Agent 负责意图识别、规划执行、RAG 检索、记忆管理与流式交互。\\
\textbf{项目职责}：
\begin{enumerate}[parsep=0.5ex]
    \item \textbf{LangGraph 编排与路由}：基于 LangGraph 构建状态图主链路，统一编排语义缓存、Query Rewrite、意图路由、上下文桥接、动作规划、工具执行与最终响应；设计 ACTION / TICKET / RAG / CHITCHAT 多专家意图打分，结合最高分、分差与置信度阈值完成分层路由决策。
    \item \textbf{Tool Calling 与 MCP 扩展}：设计标准 Tool Schema，将查票、路线查询、订票、订单查询、退票等能力封装为工具；接入外部天气 MCP 工具，支持 Agent 联动查询车次与出发/到达城市天气，生成“车次可用性 + 出行风险提醒”的多工具综合建议。
    \item \textbf{Plan-and-Execute 与安全执行}：针对购票、退票等高风险动作请求引入 Plan-and-Execute、显式确认态、执行失败兜底、Replan 与副作用工具去重机制，避免误触发和重复订票/退票，并在 trace 中记录规划、确认、执行与重规划结果。
    \item \textbf{文档入库与混合检索}：实现 PDF/DOCX/Markdown/TXT 规章文档解析、条款级切分、hash 去重与入库状态追踪，支持写入 OpenSearch BM25 与 Milvus 向量索引；结合 RRF、Reranker、Planned RAG、Redis 语义缓存与长短期记忆优化复杂政策问答。
    \item \textbf{Agent Eval 与质量评估}：构建 Agent golden cases，核心用例全量通过，覆盖工具选择、参数正确性、MCP trace、副作用安全与坏 case 回归；接入 DeepEval G-Eval + 智谱 judge 评估回答质量，RAGAS 历史结果达到 context precision 0.881、context recall 0.935。
\end{enumerate}

\datedsubsection{\textbf{12306 高并发售票核心网关}（交易内核/后端开发）}{2026.03 - 2026.06}
\textbf{项目技术}：Spring Boot, Redis, Redisson, Caffeine, RBloomFilter, RocketMQ, Sentinel, MyBatis-Plus\\
\textbf{项目描述}：作为 Python Agent 的票务交易工具内核，面向 12306 高并发售票场景，围绕查票、订票、退款等核心链路完成缓存治理、库存预扣、事务一致性与补偿闭环设计。\\
\textbf{项目职责}：
\begin{enumerate}[parsep=0.5ex]
    \item \textbf{缓存与高并发治理}：基于 Caffeine + Redis + RBloomFilter 构建多级缓存，结合 Redisson 分布式锁与基于 CompletableFuture 的 singleflight 回源合并治理热点 Key；在查票热路径压测中，800 并发下达 5409.54 QPS、0\% HTTP 错误率，P95 344ms。
    \item \textbf{库存扣减与分桶模型}：基于 Redis Lua 实现库存原子预扣减，设计 Segmented Stock 分桶库存模型分散热点流量，降低高并发订票场景下的库存争抢与超卖风险。
    \item \textbf{订单状态机与事务一致性}：设计待支付、已支付、已取消、退款中、已退款等订单状态流转，结合状态条件更新与业务幂等号处理支付回调、超时关单、MQ 重试和重复退票；通过 RocketMQ 事务消息串联 Redis 预扣、订单落库与 MySQL 扣减，在 200 用户订票压测中达 510.94 TPS、0\% 错误率，未观察到超卖。
    \item \textbf{退款补偿与稳定性保护}：设计 refund 字段与 Redis Lua 幂等回补机制，解决 MQ 重试导致的重复补库存问题；基于 RocketMQ 延迟消息实现超时关单，并结合 Sentinel 完成限流与热点保护，为 Agent 高风险工具调用提供可靠后端边界。
\end{enumerate}

\end{document}
